%% =====================================================================
%% Canovaccio della discussione NDS (prof. Baggio), 26/11/2026, ore 9:30
%% Solo il discorso: per ogni slide le frasi da dire, nell'ordine della slide.
%% =====================================================================
%% COME SI AGGIORNA
%% 1. Ogni slide e' un blocco
%%       \begin{slide}{<label del frame>}{<titolo>}{<tempo>}
%%         \begin{discorso}
%%           \item prima frase
%%           \item seconda frase
%%         \end{discorso}
%%       \end{slide}
%%    La <label> e' quella del frame in main.tex, es. [label=m-bridge].
%%    Frame a due passi: \begin{slidedue}{...}{...}{...}; \clic segna il passo 2.
%%    Pagine senza label (titolo, outline, Thank you): \begin{slidepag}{<pagina>}{...}{...}.
%% 2. Miniature, numeri di slide e di pagina NON si scrivono a mano: vengono
%%    da ../main.pdf e ../main.aux. Se cambi la presentazione, ricompila prima
%%    main.tex e poi questo file (pdflatex canovaccio.tex, due volte).
%% 3. Nel testo: \fr{label} = numero della slide; \pg{label} = pagina del PDF.
%% 4. Regole del discorso: niente numeri di teoremi, corollari o casi (si dice il
%%    contenuto o il nome); si dice solo cio' che e' sulla slide; termini tecnici in
%%    inglese come sulle slide (consensus, disagreement, prejudice, susceptibility,
%%    attachment, augmented, social power, influence centrality, ...).
%% 5. Domande probabili, termini da evitare e cose nostre: file domande.tex.
%% =====================================================================
\documentclass[a4paper,10pt]{article}
\usepackage[T1]{fontenc}
% babel in italiano: con TeX Live completo usa italian.ldf, altrimenti il file .ini di babel
\IfFileExists{italian.ldf}{\usepackage[italian]{babel}}%
  {\usepackage{babel}\babelprovide[import,main]{italian}}
\usepackage{lmodern}
\usepackage{microtype}
\usepackage[a4paper,margin=1.5cm,top=1.4cm,bottom=1.7cm]{geometry}
\usepackage{amsmath,amssymb}
\usepackage{graphicx}
\usepackage{xcolor}
\usepackage{enumitem}
\usepackage{refcount}
\usepackage{xr-hyper}
\usepackage[hidelinks]{hyperref}
\externaldocument[P-]{../main}  % la presentazione nuova
\externaldocument[O-]{../../presentation_v6/main}  % la presentazione precedente
\pdfsuppresswarningpagegroup=1

% ---------------------------------------------------------------------
% Riferimenti alla presentazione
\newcommand{\presentazione}{../main.pdf}
\newcommand{\pg}[1]{\getpagerefnumber{P-#1}}  % pagina del PDF
\newcommand{\fr}[1]{\getrefnumber{P-#1}}      % numero della slide
\newcommand{\uno}{\mathbf{1}}

% ---------------------------------------------------------------------
% Impaginazione
\setlength{\parindent}{0pt}
\setlength{\parskip}{2pt}
\raggedbottom
\definecolor{unipd}{RGB}{155,0,20}

\newcommand{\mini}[1]{{\setlength{\fboxsep}{0pt}%
  \fbox{\includegraphics[page=#1,width=\dimexpr\linewidth-2\fboxrule\relax]{\presentazione}}}}
\newcommand{\testata}[4]{% numero, pagine, titolo, tempo
  {\color{unipd}\bfseries Slide #1 \textperiodcentered\ #3}\hfill{\small\itshape #4}\par
  {\scriptsize\color{gray} pagina #2 del PDF}\par\smallskip}

% Blocco di una slide: miniatura a sinistra, discorso a destra
\newenvironment{slide}[3]{%
  \par\addvspace{12pt}\noindent
  \begin{minipage}[t]{0.36\textwidth}\vspace{0pt}\mini{\pg{#1}}\end{minipage}\hfill
  \begin{minipage}[t]{0.61\textwidth}\vspace{0pt}%
  \testata{\fr{#1}}{\pg{#1}}{#2}{#3}}%
 {\end{minipage}\par}
% Frame a due passi: le due miniature una sotto l'altra
\newenvironment{slidedue}[3]{%
  \par\addvspace{12pt}\noindent
  \begin{minipage}[t]{0.36\textwidth}\vspace{0pt}\mini{\pg{#1<1>}}\par\smallskip
    \mini{\pg{#1<2>}}\end{minipage}\hfill
  \begin{minipage}[t]{0.61\textwidth}\vspace{0pt}%
  \testata{\fr{#1}}{\pg{#1<1>}--\pg{#1<2>}}{#2}{#3}}%
 {\end{minipage}\par}
% Pagine senza label: si indica la pagina del PDF
\newenvironment{slidepag}[3]{%
  \par\addvspace{12pt}\noindent
  \begin{minipage}[t]{0.36\textwidth}\vspace{0pt}\mini{#1}\end{minipage}\hfill
  \begin{minipage}[t]{0.61\textwidth}\vspace{0pt}%
  {\color{unipd}\bfseries #2}\hfill{\small\itshape #3}\par
  {\scriptsize\color{gray} pagina #1 del PDF}\par\smallskip}%
 {\end{minipage}\par}

% Il discorso: una frase (o due) per punto
\newenvironment{discorso}{\begin{itemize}[nosep,leftmargin=1em,label={\color{gray}\textbullet},itemsep=2pt]}{\end{itemize}}
\newcommand{\clic}{{\color{unipd}\scriptsize[passo 2]}\ }
\newcommand{\parte}[2]{\section*{#1\hfill{\normalsize\itshape #2}}}


% testo tolto: barrato parola per parola (senza pacchetti esterni)
\newsavebox{\barratobox}
\newcommand{\barrato}[1]{\leavevmode\sbox{\barratobox}{#1}%
  \rlap{\rule[0.55ex]{\wd\barratobox}{0.5pt}}\usebox{\barratobox}}
\newcommand{\vecchia}{../../presentation_v6/main.pdf}
\newcommand{\pgO}[1]{\getpagerefnumber{O-#1}}
\newcommand{\frO}[1]{\getrefnumber{O-#1}}
\newcommand{\miniO}[1]{{\setlength{\fboxsep}{0pt}%
  \fbox{\includegraphics[page=#1,width=\dimexpr\linewidth-2\fboxrule\relax]{\vecchia}}}}
\newcommand{\etichetta}[2]{\,{\setlength{\fboxsep}{1.5pt}\colorbox{#1}{\scriptsize\bfseries\color{white}#2}}}
\newcommand{\nuova}{\etichetta{blue}{NUOVA}}
\newcommand{\spostata}{\etichetta{orange}{SPOSTATA}}
\newcommand{\divisa}{\etichetta{teal}{DIVISA}}
\newcommand{\toltatesto}[1]{{\color{red}\barrato{#1}}}
\begin{document}
\begin{center}
  {\LARGE\bfseries Preview delle modifiche}\\[3pt]
  {\large Presentazione e canovaccio, discussione NDS del 26/11}\\[2pt]
  {\small\color{gray} Versione del \today. Confronto con la versione precedente: canovaccio del 17/11 (prima di queste note) e presentazione del 16/11.}
\end{center}

\textbf{Legenda.} {\color{blue}Blu}: testo nuovo. {\color{red}\barrato{Rosso barrato}}: testo
tolto (le formule tolte sono solo in rosso). \nuova: slide o blocco nuovo.
\spostata: slide spostata. \divisa: slide divisa in due.
\section*{Slide modificate: prima e dopo}
Solo le slide cambiate, nel nuovo ordine. La parte FJ diventa: The FJ model, Convergence in the FJ model, Convergence with P-independent agents, Final opinions and the augmented graph, Back to the question, Influence centrality, FDG and FJ on the same networks, Stubborn agent: state or input, poi gli esempi. In tutto 30 slide (prima 28).
\par\addvspace{14pt}\noindent\begin{minipage}{\textwidth}
{\color{unipd}\bfseries Slide \fr{m-fdg-why} \textperiodcentered\ Why the same social power}\par\smallskip
{\small Tolta la colonna «Check (report, Section 2.2.3)». Resta il bilancio, detto per intero: se tutti passano la stessa quantità $b$, tutti raccolgono $b$; quindi $w_j = b$ per gli agenti 2--6 e $w_1 = t\,w_1 + b$ per l'agente 1. L'algebra resta nel backup «Derivation», dal pulsante della slide precedente.}\par\smallskip
\noindent\begin{minipage}[t]{0.36\textwidth}\vspace{0pt}{\scriptsize\color{gray} Prima}\par
\miniO{\pgO{m-fdg-why}}\par\smallskip
\end{minipage}\hfill\begin{minipage}[t]{0.61\textwidth}\vspace{0pt}{\scriptsize\color{gray} Dopo}\par
\mini{\pg{m-fdg-why}}\par\smallskip
\end{minipage}\par
\end{minipage}\par
\par\addvspace{14pt}\noindent\begin{minipage}{\textwidth}
{\color{unipd}\bfseries Slide \fr{m-fj-conv}--\fr{m-fj-conv2} \textperiodcentered\ Convergence in the FJ model}\divisa\par\smallskip
{\small Divisa in due. La prima: definizioni, P-dependent, condizione di stabilità asintotica e l'intuizione scritta sulla slide, $e(k+1) = \Lambda W\, e(k)$. La seconda: il caso generale con agenti P-independent, con il significato di $W^{22}$ scritto sulla slide.}\par\smallskip
\noindent\begin{minipage}[t]{0.36\textwidth}\vspace{0pt}{\scriptsize\color{gray} Prima}\par
\miniO{\pgO{m-fj-conv}}\par\smallskip
\end{minipage}\hfill\begin{minipage}[t]{0.61\textwidth}\vspace{0pt}{\scriptsize\color{gray} Dopo}\par
\mini{\pg{m-fj-conv}}\par\smallskip
\mini{\pg{m-fj-conv2}}\par\smallskip
\end{minipage}\par
\end{minipage}\par
\par\addvspace{14pt}\noindent\begin{minipage}{\textwidth}
{\color{unipd}\bfseries Slide \fr{m-back-question} \textperiodcentered\ Back to the question}\spostata\par\smallskip
{\small Spostata subito dopo «Final opinions and the augmented graph». In alto la domanda ripetuta; la risposta parte dal modello augmented letto al contrario. Due passi: prima la riscrittura, poi la risposta.}\par\smallskip
\noindent\begin{minipage}[t]{0.36\textwidth}\vspace{0pt}{\scriptsize\color{gray} Prima}\par
\miniO{\pgO{m-back-question}}\par\smallskip
\end{minipage}\hfill\begin{minipage}[t]{0.61\textwidth}\vspace{0pt}{\scriptsize\color{gray} Dopo}\par
\mini{\pg{m-back-question<1>}}\par\smallskip
\mini{\pg{m-back-question<2>}}\par\smallskip
\end{minipage}\par
\end{minipage}\par
\par\addvspace{14pt}\noindent\begin{minipage}{\textwidth}
{\color{unipd}\bfseries Slide \fr{m-fj-centrality} \textperiodcentered\ Influence centrality}\par\smallskip
{\small Resta solo la definizione: il caso $\alpha = 0$ e il Lemma 3.5 passano su «Centrality meets social power», dove si usano.}\par\smallskip
\noindent\begin{minipage}[t]{0.36\textwidth}\vspace{0pt}{\scriptsize\color{gray} Prima}\par
\miniO{\pgO{m-fj-centrality}}\par\smallskip
\end{minipage}\hfill\begin{minipage}[t]{0.61\textwidth}\vspace{0pt}{\scriptsize\color{gray} Dopo}\par
\mini{\pg{m-fj-centrality}}\par\smallskip
\end{minipage}\par
\end{minipage}\par
\par\addvspace{14pt}\noindent\begin{minipage}{\textwidth}
{\color{unipd}\bfseries Slide \fr{m-fj-networks} \textperiodcentered\ FDG and FJ on the same networks}\nuova\par\smallskip
{\small Prima degli esempi: su quali reti lavoriamo ($W = A(0)$, $u = x(0)$) e come si mette lo stesso peso nei due modelli, in tre righe (agente 1 stubborn; solo l'agente 1 con peso $t$; tutti con peso 0.2). Nella seconda riga le due equazioni differiscono solo per $x_1(k)$ al posto di $u_1$.}\par\smallskip
\noindent\begin{minipage}[t]{0.36\textwidth}\vspace{0pt}{\scriptsize\color{gray} Prima}\par
{\small ---}
\end{minipage}\hfill\begin{minipage}[t]{0.61\textwidth}\vspace{0pt}{\scriptsize\color{gray} Dopo}\par
\mini{\pg{m-fj-networks}}\par\smallskip
\end{minipage}\par
\end{minipage}\par
\par\addvspace{14pt}\noindent\begin{minipage}{\textwidth}
{\color{unipd}\bfseries Slide \fr{m-all-prej} \textperiodcentered\ Every agent prejudiced}\par\smallskip
{\small Solo la didascalia: lo stesso attachment per tutti, prejudice diversi.}\par\smallskip
\noindent\begin{minipage}[t]{0.36\textwidth}\vspace{0pt}{\scriptsize\color{gray} Prima}\par
\miniO{\pgO{m-all-prej}}\par\smallskip
\end{minipage}\hfill\begin{minipage}[t]{0.61\textwidth}\vspace{0pt}{\scriptsize\color{gray} Dopo}\par
\mini{\pg{m-all-prej}}\par\smallskip
\end{minipage}\par
\end{minipage}\par
\par\addvspace{14pt}\noindent\begin{minipage}{\textwidth}
{\color{unipd}\bfseries Slide \fr{m-aux} \textperiodcentered\ Where no prejudice arrives}\par\smallskip
{\small Passo 1: la rete e tre esempi, ognuno con cosa ci aspettiamo. Passo 2: i risultati, con lo stesso nome degli esempi.}\par\smallskip
\noindent\begin{minipage}[t]{0.36\textwidth}\vspace{0pt}{\scriptsize\color{gray} Prima}\par
\miniO{\pgO{m-aux<1>}}\par\smallskip
\miniO{\pgO{m-aux<2>}}\par\smallskip
\end{minipage}\hfill\begin{minipage}[t]{0.61\textwidth}\vspace{0pt}{\scriptsize\color{gray} Dopo}\par
\mini{\pg{m-aux<1>}}\par\smallskip
\mini{\pg{m-aux<2>}}\par\smallskip
\end{minipage}\par
\end{minipage}\par
\par\addvspace{14pt}\noindent\begin{minipage}{\textwidth}
{\color{unipd}\bfseries Slide \fr{m-lemma} \textperiodcentered\ Centrality meets social power}\par\smallskip
{\small Due passi. Passo 1: le due misure, perché si confrontano in parallelo, il Lemma 3.5 (spostato qui) e cosa ci aspettiamo. Passo 2: il grafico.}\par\smallskip
\noindent\begin{minipage}[t]{0.36\textwidth}\vspace{0pt}{\scriptsize\color{gray} Prima}\par
\miniO{\pgO{m-lemma}}\par\smallskip
\end{minipage}\hfill\begin{minipage}[t]{0.61\textwidth}\vspace{0pt}{\scriptsize\color{gray} Dopo}\par
\mini{\pg{m-lemma<1>}}\par\smallskip
\mini{\pg{m-lemma<2>}}\par\smallskip
\end{minipage}\par
\end{minipage}\par
\par\smallskip\noindent Backup: il pulsante di ritorno di «Convergence and spectrum» porta ora a «Convergence with P-independent agents», quello di «Influence centrality» a «Centrality meets social power». Gli enunciati sono gli stessi della versione precedente (il Teorema 3.2 cambia solo l'impaginazione).\par
\clearpage
\section*{Canovaccio: modifiche rispetto alla versione precedente}
Il canovaccio nuovo segue il nuovo ordine delle slide. Regole applicate ovunque: niente numeri di teoremi, corollari o casi (si dice il contenuto o il nome); si dice solo ciò che è sulla slide; termini in inglese come sulle slide (consensus, disagreement, prejudice, susceptibility, attachment, augmented).
\parte{Parte 1 — Introduzione}{3:00}
\begin{slidepag}{1}{Titolo}{0:15}
  \begin{discorso}
    \item Buongiorno. Presento il mio progetto su stubbornness e social power nella dinamica delle opinioni.
    \item Studio due modelli, il modello di French–Harary–DeGroot e il modello di Friedkin–Johnsen: come cambia il loro comportamento quando gli agenti resistono agli altri, e quanto pesa ciascun agente nelle opinioni finali.
  \end{discorso}
\end{slidepag}

\begin{slidepag}{2}{Slide 2 \textperiodcentered\ Outline}{0:15}
  \begin{discorso}
    \item La presentazione ha quattro parti: introduzione, modello FDG, modello FJ, conclusioni.
    \item Alcune slide hanno dei pulsanti in basso a destra: aprono le slide di backup, con definizioni e derivazioni. Le uso solo se vuole entrare nel dettaglio.
  \end{discorso}
\end{slidepag}

\begin{slide}{m-question}{Stubbornness}{1:00}
  \begin{discorso}
    \item Una rete sociale è un digrafo pesato: gli agenti sono i nodi, gli archi portano l'influenza, e ogni agente aggiorna la sua opinione facendo la media delle opinioni che ascolta.
    \item Un agente è stubborn se la sua opinione resta invariata indipendentemente dalle opinioni degli altri: è la definizione del tutorial di Proskurnikov e Tempo.
    \item È il caso estremo di due modi graduati di resistere agli altri: dare peso alla propria opinione attuale, oppure restare attaccati a un'opinione fissa propria, un {\color{red}\barrato{pregiudizio.}} {\color{blue}prejudice.}
    \item Il primo è il modo del modello FDG, il secondo del modello FJ: li introduco tra poco.
    \item A entrambi i modelli, sulle stesse reti, faccio tre domande. Le opinioni convergono? Convergono alla stessa opinione per tutti, il {\color{red}\barrato{consenso,}} {\color{blue}consensus,} o a opinioni diverse, il {\color{red}\barrato{disaccordo?}} {\color{blue}disagreement?} E quanto pesa l'opinione iniziale di ciascun agente in quelle finali?
    \item Prima di tutto, da dove vengono i due modelli.
  \end{discorso}
\end{slide}

\begin{slide}{m-models}{Models of opinion dynamics}{0:40}
  \begin{discorso}
    \item Il quadro storico è quello del tutorial di Proskurnikov e Tempo, la fonte principale del lavoro, insieme alle slide del corso e al libro di Bullo.
    \item A sinistra i modelli che fanno solo medie: in tempo discreto il modello di French, Harary e DeGroot; in tempo continuo quello di Abelson.
    \item Per i gruppi reali il {\color{red}\barrato{consenso}} {\color{blue}consensus} è una previsione troppo forte: spesso restano divisi.
    \item Per questo il modello di Friedkin e Johnsen, negli anni Novanta, aggiunge alle medie opinioni esterne costanti; è stato validato sperimentalmente su gruppi piccoli e medi. In tempo continuo il corrispondente è il modello di Taylor.
    \item In questo lavoro tratto il modello FDG, visto nel corso, e il modello FJ, che è nuovo; Abelson e Taylor restano solo nella panoramica.
    \item Richiamo solo le nozioni del corso che mi servono.
  \end{discorso}
\end{slide}

\begin{slide}{m-prelim}{Preliminaries}{0:50}
  \begin{discorso}
    \item Sono nozioni del corso, quindi solo i punti che userò.
    \item Un arco da $i$ a $j$, con peso $a_{ij}$, vuol dire che l'agente $i$ è influenzato dall'agente $j$: è la convenzione del corso; il tutorial disegna gli archi al contrario.
    \item Un digrafo fortemente connesso è periodico se un intero maggiore di uno divide la lunghezza di ogni ciclo semplice; un self-loop lo rende aperiodico.
    \item Nel condensation digraph ogni componente fortemente connessa diventa un nodo; un sink è un gruppo influenzato solo da sé stesso. Saranno i sink a decidere l'esito.
    \item Per le matrici: semi-convergente vuol dire che il limite delle potenze esiste; convergente, o Schur stabile, che il limite è zero, cioè raggio spettrale minore di uno.
    \item Infine: irriducibile equivale a fortemente connesso, primitiva a fortemente connesso e aperiodico.
    \item Ora il primo modello.
  \end{discorso}
\end{slide}

\parte{Parte 2 — Il modello FDG}{7:00}
\begin{slide}{m-fdg-model}{The FDG model}{0:50}
  \begin{discorso}
    \item Nel modello FDG ogni agente sostituisce la sua opinione con una media pesata delle opinioni che ascolta.
    \item In forma vettoriale $x(k+1)=Ax(k)$, con $A$ la matrice di adiacenza pesata, row-stochastic: è un sistema di media a tempo discreto.
    \item Iterando, $x(k)=A^kx(0)$: il comportamento nel lungo periodo dipende dalle potenze di $A$. Una matrice row-stochastic ha sempre l'autovalore 1, e raggio spettrale 1.
    \item {\color{red}\barrato{Il}} {\color{blue}A destra la definizione: il} modello è convergente se, per ogni condizione iniziale, il limite esiste; raggiunge un {\color{red}\barrato{consenso}} {\color{blue}consensus} se, per ogni condizione iniziale, le opinioni finali sono tutte uguali.
    \item Gli esiti possibili sono quindi tre: nessuna convergenza, cioè $A$ non semi-convergente; {\color{red}\barrato{consenso;}} {\color{blue}consensus;} e {\color{red}\barrato{disaccordo,}} {\color{blue}disagreement,} cioè il modello converge ma per qualche condizione iniziale le opinioni finali non sono tutte uguali.
    \item Quale dei tre si ha si legge in tre modi equivalenti.
  \end{discorso}
\end{slide}

\begin{slide}{m-fdg-rosetta}{Which outcome: the Rosetta stone}{0:40}
  \begin{discorso}
    \item L'esito si legge sullo spettro di $A$, sulla matrice o sul digrafo: la Rosetta stone del corso collega i tre punti di {\color{red}\barrato{vista.}} {\color{blue}vista, e ogni riga della tabella è un caso.}
    \item Sul digrafo conta solo il condensation digraph: quanti sink {\color{red}\barrato{ha}} {\color{blue}ha,} e se sono periodici.
    \item {\color{red}\barrato{Un}} {\color{blue}Se c'è un} solo sink, {\color{blue}ed è} aperiodico, {\color{red}\barrato{dà}} {\color{blue}si ha consensus: è} il {\color{red}\barrato{consenso:} \barrato{casi} \barrato{(i)}} {\color{blue}caso di una rete fortemente connessa} e {\color{red}\barrato{(iii).} \barrato{Almeno} \barrato{due} \barrato{sink,} \barrato{tutti} \barrato{aperiodici,} \barrato{danno} \barrato{il} \barrato{disaccordo:} \barrato{caso} \barrato{(iv).} \barrato{Un}} {\color{blue}aperiodica, o di un unico} sink {\color{red}\barrato{periodico} \barrato{impedisce} \barrato{la} \barrato{convergenza:} \barrato{casi} \barrato{(ii)} \barrato{e} \barrato{(v).}} {\color{blue}aperiodico.}
    \item {\color{blue}Se ci sono almeno due sink, tutti aperiodici, si ha disagreement. Se almeno un sink è periodico, il modello non converge.}
    \item {\color{red}\barrato{Mi}} {\color{blue}In basso il fatto che mi} servirà {\color{red}\barrato{soprattutto} \barrato{il} \barrato{Corollario} \barrato{1.12:}} {\color{blue}di più:} se il digrafo è fortemente connesso e ha almeno un self-loop, $A$ è primitiva.
    \item Il lavoro si concentra sul {\color{red}\barrato{consenso,}} {\color{blue}consensus,} perché è il caso in cui il peso di ciascun agente si può misurare. Al {\color{red}\barrato{disaccordo}} {\color{blue}disagreement} dedico la slide dopo.
  \end{discorso}
\end{slide}

\begin{slide}{m-fdg-consensus}{Consensus and social power}{0:50}
  \begin{discorso}
    \item Il teorema del {\color{red}\barrato{consenso}} {\color{blue}consensus} del corso dà tre condizioni equivalenti.
    \item Sullo spettro: l'autovalore 1 è semplice e gli altri hanno modulo minore di uno. Sulla matrice: le potenze di $A$ convergono a $\uno w^\top$. Sul grafo: esiste un nodo globally reachable, e il sottografo dei nodi globally reachable è aperiodico.
    \item In questo caso $w$ è l'autovettore sinistro dominante di $A$, con $w_i$ positivo se e solo se il nodo $i$ è globally reachable, e il valore di {\color{red}\barrato{consenso}} {\color{blue}consensus} è $w^\top x(0)$.
    \item {\color{blue}A destra:} $w_i$ è la social power dell'agente {\color{red}$i$:} {\color{blue}$i$,} il peso della sua opinione iniziale nel valore di {\color{red}\barrato{consenso.}} {\color{blue}consensus.}
    \item Richiede il {\color{red}\barrato{consenso,}} {\color{blue}consensus,} perché solo allora le righe del limite di $A^k$ sono tutte uguali a $w^\top$.
    \item {\color{red}\barrato{Se}} {\color{blue}Sotto, l'average consensus: se} $A$ è doubly-stochastic e vale la condizione sul grafo, $w$ è uniforme e il {\color{red}\barrato{consenso}} {\color{blue}consensus} è la media delle opinioni iniziali.
    \item Nei termini del tutorial: il modello converge se e solo se $A$ è regular, e raggiunge il {\color{red}\barrato{consenso}} {\color{blue}consensus} se e solo se $A$ è fully regular.
    \item E se il {\color{red}\barrato{consenso}} {\color{blue}consensus} non c'è?
  \end{discorso}
\end{slide}

\begin{slide}{m-fdg-disagreement}{Disagreement and stubbornness}{0:40}
  \begin{discorso}
    \item Se il condensation digraph ha almeno due sink, il teorema del {\color{red}\barrato{disaccordo}} {\color{blue}disagreement} dà tre condizioni equivalenti: l'autovalore 1 è semi-semplice, con molteplicità pari al numero di sink; $A$ è semi-convergente; ogni sink è aperiodico.
    \item Allora ogni sink raggiunge un {\color{red}\barrato{consenso}} {\color{blue}consensus} al suo interno, e ogni agente fuori dai sink converge a una combinazione convessa dei valori dei sink.
    \item Le righe del limite di $A^k$ sono diverse: nessun vettore fa la parte di $w$, e la social power non è definita.
    \item A destra la definizione di stubborn del tutorial. Nel modello FDG l'agente $i$ è stubborn se e solo se $a_{ii}=1$: è un sink fatto di un solo nodo, aperiodico grazie al self-loop.
    \item Se tutti gli agenti sono influenzati da lui, direttamente o attraverso altri, è l'unico nodo globally reachable, e la sua social power è 1.
    \item Ora le due reti su cui confronto i due modelli.
  \end{discorso}
\end{slide}

\begin{slide}{m-networks}{Two networks, one agent}{0:40}
  \begin{discorso}
    \item Uso due reti di sei agenti: un ciclo diretto, in cui ogni agente ascolta solo il successivo, e un anello bidirezionale, in cui ogni agente ascolta i due vicini con peso un mezzo.
    \item Per $t=0$ entrambe sono fortemente connesse e periodiche, con periodo 6 il ciclo e 2 l'anello, e doubly-stochastic: l'autovettore sinistro dell'autovalore 1 è uniforme, nessun agente è favorito.
    \item Cambio solo la riga dell'agente 1: $t$ è il {\color{blue}self-weight dell'agente 1, il} peso che {\color{red}\barrato{l'agente} \barrato{1}} dà alla propria opinione attuale. Si va dalle reti originali, $t=0$, all'agente 1 stubborn, $t=1$.
    \item La tabella in basso: per $t$ fra 0 e 1 le reti sono fortemente connesse e aperiodiche, e non più doubly-stochastic; per $t=1$ non sono più fortemente connesse, perché l'agente 1 non ascolta nessun altro.
    \item La condizione iniziale è sempre la stessa: opinioni da 0 a 1, con passo 0.2.
    \item Vediamo cosa succede alle opinioni.
  \end{discorso}
\end{slide}

\begin{slide}{m-fdg-conv}{One self-weight brings agreement}{0:50}
  \begin{discorso}
    \item Le colonne sono $t=0$, un mezzo e 1; sopra il ciclo, sotto l'anello.
    \item Per $t=0$ le reti sono periodiche e le opinioni oscillano per sempre: nessuna {\color{red}\barrato{convergenza,} \barrato{caso} \barrato{(ii).}} {\color{blue}convergenza.} Sullo spettro ci sono autovalori diversi da 1 sul cerchio unitario.
    \item Basta un self-weight $t$ positivo, di un solo agente, e il gruppo raggiunge il {\color{red}\barrato{consenso.}} {\color{blue}consensus.}
    \item Per $t$ fra 0 e 1 la rete è fortemente connessa {\color{red}\barrato{con}} {\color{blue}e ha} un self-loop, quindi $A(t)$ è {\color{red}\barrato{primitiva} \barrato{per} \barrato{il} \barrato{Corollario} \barrato{1.12:} \barrato{caso} \barrato{(i).}} {\color{blue}primitiva: consensus.}
    \item Per $t=1$ l'agente 1 è stubborn ed è l'unico nodo globally reachable: tutti finiscono sulla sua opinione iniziale, zero. {\color{red}\barrato{Caso} \barrato{(iii).}}
    \item Per ogni $t$ positivo gli autovalori diversi da 1 sono tutti strettamente dentro il disco unitario.
    \item Questo è il primo punto del lavoro: la testardaggine ribalta la dinamica. Un agente che dà peso alla propria opinione attuale, senza essere stubborn, basta perché su una rete fortemente connessa tutto il gruppo si metta d'accordo.
    \item Con il {\color{red}\barrato{consenso}} {\color{blue}consensus} possiamo misurare il peso di ciascun agente.
  \end{discorso}
\end{slide}

\begin{slide}{m-fdg-sp}{Social power grows gradually}{0:45}
  \begin{discorso}
    \item Per $t$ in $(0,1]$ c'è {\color{red}\barrato{il} \barrato{consenso,}} {\color{blue}consensus,} e la social power ha la forma chiusa a destra, la stessa sulle due reti.
    \item In alto: la social power dell'agente 1 cresce gradualmente con il suo self-weight, da $1/6$, che è il limite per $t$ che tende a 0, fino a 1 per $t=1$. A $t=\tfrac12$ vale $2/7$, e ogni altro agente $1/7$.
    \item In basso il valore di {\color{red}\barrato{consenso,}} {\color{blue}consensus,} $3(1-t)/(6-5t)$: da 0.5 nel limite per $t$ che tende a 0, fino a $x_1(0)=0$; a $t=\tfrac12$ vale $3/7$.
    \item I valori sono gli stessi sulle due reti: la rete decide come le opinioni convergono, non dove.
    \item In termini sociali: l'inerzia guadagna influenza poco alla volta, e per $t$ fra 0 e 1 ogni opinione iniziale entra ancora nel {\color{red}\barrato{consenso.}} {\color{blue}consensus.}
    \item Perché la stessa social power su due reti così diverse?
  \end{discorso}
\end{slide}

\begin{slide}{m-fdg-why}{Why the same social power}{0:50 {\color{red}\barrato{0:45}}}
  \begin{discorso}
    \item {\color{blue}Per $t$ in $(0,1]$ c'è consensus, e la social power è l'unica soluzione di $w^\top A(t)=w^\top$ con somma uno.}
    \item {\color{red}\barrato{La} \barrato{social} \barrato{power} \barrato{risolve}} {\color{red}$w^\top A(t)=w^\top$,} {\color{red}\barrato{che} \barrato{si} \barrato{legge}} {\color{blue}A sinistra la leggo} come un bilancio: ogni agente passa la sua social power agli agenti che ascolta, in proporzione ai pesi, e $w_j$ è quello che l'agente $j$ {\color{red}\barrato{raccoglie.}} {\color{blue}raccoglie, compresa la parte che trattiene.}
    \item {\color{red}$A(0)$} {\color{red}\barrato{è} \barrato{doubly-stochastic:}} {\color{blue}Nella rete originale, $A(0)$,} ogni agente è ascoltato con peso totale {\color{red}\barrato{1,} \barrato{quindi} \barrato{il} \barrato{bilancio} \barrato{torna} \barrato{quando} \barrato{ogni} \barrato{agente} \barrato{passa} \barrato{la} \barrato{stessa} \barrato{quantità.}} {\color{blue}uno: è doubly-stochastic.}
    \item {\color{red}\barrato{Gli} \barrato{agenti} \barrato{da} \barrato{2} \barrato{a} \barrato{6} \barrato{non} \barrato{hanno} \barrato{self-loop} \barrato{e} \barrato{passano} \barrato{tutto} \barrato{il} \barrato{loro}} {\color{red}$b$;} {\color{red}\barrato{l'agente} \barrato{1} \barrato{trattiene} \barrato{la} \barrato{quota}} {\color{red}$t$} {\color{red}\barrato{e} \barrato{passa}} {\color{red}$(1-t)w_1$.}
    \item {\color{red}\barrato{Quindi}} {\color{red}$(1-t)w_1=b$:} {\color{red}\barrato{l'agente} \barrato{1} \barrato{pesa}} {\color{red}$1/(1-t)$} {\color{red}\barrato{volte} \barrato{ogni} \barrato{altro} \barrato{agente,}} {\color{red}$2/7$} {\color{red}\barrato{contro}} {\color{red}$1/7$} {\color{red}\barrato{a}} {\color{red}$t=\tfrac12$.}
    \item {\color{red}\barrato{Conta}} {\color{blue}Cambiando $t$, l'agente 1 cambia} solo quanto {\color{red}\barrato{trattiene} \barrato{l'agente} \barrato{1,} \barrato{non} \barrato{chi} \barrato{ascolta:} \barrato{per} \barrato{questo}} {\color{red}$w(t)$} {\color{red}\barrato{è} \barrato{la} \barrato{stessa} \barrato{sulle} \barrato{due} \barrato{reti.}} {\color{blue}trattiene; il resto lo passa con i pesi della rete originale.}
    \item {\color{blue}Quindi, se ogni agente passa la stessa quantità $b$, ognuno raccoglie esattamente $b$.}
    \item {\color{blue}A destra la conseguenza. Gli agenti da 2 a 6 non trattengono niente: la loro social power è $b$. L'agente 1 trattiene $t\,w_1$, quindi $w_1=t\,w_1+b$, cioè $b=(1-t)\,w_1$.}
    \item {\color{blue}Con la somma uguale a uno si ritrova la formula della slide precedente, e siccome la soluzione è unica, è proprio la social power. L'agente 1 pesa $1/(1-t)$ volte ogni altro agente: $2/7$ contro $1/7$ a $t=\tfrac12$.}
    \item {\color{blue}In basso: conta solo quanto trattiene l'agente 1, non chi ascolta; per questo $w(t)$ è la stessa sulle due reti.}
    \item {\color{red}\barrato{A} \barrato{destra} \barrato{la} \barrato{verifica} \barrato{del} \barrato{report:} \barrato{sostituendo}} {\color{red}$w=(w_1,b,\dots,b)$} {\color{red}\barrato{si} \barrato{ritrova}} {\color{red}$b=(1-t)w_1$.}
    \item Per $t$ che tende a 0, $w$ tende al vettore uniforme; ma a $t=0$ non c'è {\color{red}\barrato{consenso,}} {\color{blue}consensus,} e quel vettore è solo l'autovettore sinistro di $A(0)$.
    \item Finora al massimo un agente stubborn. Vediamo cosa succede con due.
  \end{discorso}
\end{slide}

\begin{slide}{m-bridge}{Two stubborn agents}{1:00}
  \begin{discorso}
    \item Come {\color{red}\barrato{nella} \barrato{Figura} \barrato{2}} {\color{blue}nell'esempio} con $t=1$, l'agente 1 è stubborn; ora sull'anello lo è anche l'agente 4, quello opposto.
    \item Il condensation digraph ha due sink, $\{1\}$ e $\{4\}$, aperiodici per i self-loop.
    \item Per il {\color{red}\barrato{Teorema} \barrato{2.4}} {\color{blue}teorema del disagreement} il gruppo converge a un {\color{red}\barrato{disaccordo:}} {\color{blue}disagreement:} gli agenti 2, 3, 5 e 6 finiscono su combinazioni convesse delle opinioni iniziali di 1 e di 4.
    \item Le righe del limite di $A^k$ sono diverse, quindi la social power non è definita.
    \item Qui mi fermo con il modello FDG, per due motivi: è già visto nel corso, e la social power, che è ciò che mi interessa, richiede il {\color{red}\barrato{consenso.}} {\color{blue}consensus.}
    \item Però lascio aperta una domanda. Gli agenti 1 e 4 non cambiano mai opinione: sono membri del gruppo, o sono opinioni fisse, esterne, che gli altri ascoltano?
    \item E in generale: un agente la cui opinione resta invariata indipendentemente dalle opinioni degli altri è un membro del gruppo, o un'opinione fissa esterna che gli altri ascoltano?
    \item Ci torno più avanti. Ora il secondo modello, il modello di Friedkin–Johnsen.
  \end{discorso}
\end{slide}

\parte{Parte 3 — Il modello FJ, teoria}{5:35 {\color{red}\barrato{6:00}}}
\begin{slide}{m-fj-model}{The FJ model}{1:00}
  \begin{discorso}
    \item Nel modello FJ ogni agente pesa con $\lambda_i$ la media delle opinioni che ascolta, e con il complemento $1-\lambda_i$ il suo {\color{red}\barrato{pregiudizio}} {\color{blue}prejudice} $u_i$.
    \item $W$ è row-stochastic, la matrice delle influenze sociali: ha il ruolo che $A$ aveva nel FDG.
    \item $\lambda_i$ è la {\color{red}\barrato{suscettibilità}} {\color{blue}susceptibility} dell'agente $i$ all'influenza sociale; $u$ è il vettore costante dei {\color{red}\barrato{pregiudizi,}} {\color{blue}prejudice,} opinioni fisse che gli agenti ascoltano; $1-\lambda_i$ è {\color{red}\barrato{l'attaccamento}} {\color{blue}l'attachment dell'agente} al proprio {\color{red}\barrato{pregiudizio.}} {\color{blue}prejudice.}
    \item {\color{red}\barrato{Se}} {\color{blue}A destra i casi estremi di $\lambda_i$. Con} $\lambda_i=1$ l'agente {\color{blue}non è attaccato al suo prejudice e} segue la regola del FDG; {\color{red}\barrato{con}} {\color{red}$\Lambda=I$} {\color{blue}se vale per tutti, $\Lambda=I$,} il modello è proprio il FDG con {\color{red}$A=W$,} {\color{red}\barrato{quindi}} {\color{blue}$A=W$. Per questo} da qui in poi {\color{red}$\Lambda\neq I$.} {\color{blue}almeno un agente ha $\lambda_i<1$.}
    \item {\color{blue}Un agente con $\lambda_i<1$ si dice prejudiced: il suo prejudice entra davvero nell'aggiornamento.}
    \item {\color{blue}All'estremo opposto, $\lambda_i=0$: l'agente non ascolta più nessuno, e dal primo passo in poi la sua opinione è $u_i$. È stubborn solo se $u_i$ è la sua opinione iniziale; altrimenti salta su $u_i$ al primo passo e resta lì.}
    \item {\color{red}\barrato{Se}} {\color{red}$\lambda_i<1$} {\color{red}\barrato{l'agente} \barrato{è} \barrato{prejudiced.} \barrato{Se}} {\color{red}$\lambda_i=0$} {\color{red}\barrato{l'agente} \barrato{vale}} {\color{red}$u_i$} {\color{red}\barrato{da}} {\color{red}$k=1$} {\color{red}\barrato{in} \barrato{poi,} \barrato{ma} \barrato{è} \barrato{stubborn} \barrato{solo} \barrato{se} \barrato{in} \barrato{più}} {\color{red}$u_i=x_i(0)$.}
    \item Nel {\color{red}\barrato{modello}} FJ originale {\color{red}$u=x(0)$} {\color{blue}il prejudice è l'opinione iniziale, $u=x(0)$,} e {\color{red}$1-\lambda_i=w_{ii}$;} {\color{red}\barrato{come} \barrato{nel} \barrato{tutorial,} \barrato{qui}} {\color{blue}l'attachment è il self-weight $w_{ii}$. Il tutorial} non lo {\color{red}\barrato{impongo.}} {\color{blue}impone, e io nemmeno; negli esempi userò $u=x(0)$.}
    \item {\color{red}\barrato{La} \barrato{differenza} \barrato{con}} {\color{blue}Rispetto al FDG cambia ciò a cui l'agente è legato:} il {\color{red}\barrato{FDG} \barrato{sta} \barrato{nell'oggetto:} \barrato{un}} self-weight {\color{blue}lo} lega {\color{red}\barrato{un} \barrato{agente}} alla sua opinione attuale, che {\color{red}\barrato{si} \barrato{muove} \barrato{ancora}} {\color{blue}continua a muoversi} verso gli altri; {\color{red}\barrato{un} \barrato{pregiudizio}} {\color{blue}il prejudice} lo lega a {\color{red}\barrato{un'opinione} \barrato{fissa,}} {\color{blue}un valore fisso,} che non si muove mai.
    \item {\color{red}\barrato{Torniamo} \barrato{alla} \barrato{prima} \barrato{domanda:}} {\color{blue}Ora le stesse tre domande per il modello FJ. La prima:} quando converge?
  \end{discorso}
\end{slide}

\begin{slide}{m-fj-conv}{Convergence in the FJ model}{0:50 {\color{red}\barrato{1:00}}}
  \begin{discorso}
    \item Il modello FJ converge se il limite esiste per ogni condizione {\color{red}\barrato{iniziale,} \barrato{ed}} {\color{blue}iniziale;} è asintoticamente stabile se $\Lambda W$ è convergente, cioè ha raggio spettrale minore di uno.
    \item {\color{red}\barrato{La} \barrato{nozione} \barrato{chiave} \barrato{è} \barrato{quella} \barrato{di} \barrato{agente} \barrato{P-dependent:}} {\color{blue}A sinistra:} un agente {\color{blue}è P-dependent se è} prejudiced, {\color{red}\barrato{oppure}} {\color{blue}o se è} influenzato, direttamente o attraverso altri, da un agente prejudiced. {\color{red}\barrato{Gli} \barrato{altri} \barrato{sono}} {\color{blue}Altrimenti è} P-independent.
    \item {\color{red}\barrato{Gli} \barrato{agenti} \barrato{P-independent} \barrato{non} \barrato{ascoltano} \barrato{nessun} \barrato{P-dependent:} \barrato{formano} \barrato{un} \barrato{modello} \barrato{FDG} \barrato{a} \barrato{sé,} \barrato{che} \barrato{nessun} \barrato{pregiudizio} \barrato{raggiunge.}}
    \item {\color{red}\barrato{Il} \barrato{Teorema} \barrato{3.2} \barrato{dice} \barrato{che} \barrato{il} \barrato{blocco} \barrato{dei} \barrato{P-dependent} \barrato{è} \barrato{sempre} \barrato{asintoticamente} \barrato{stabile,} \barrato{e} \barrato{che} \barrato{il} \barrato{modello} \barrato{converge} \barrato{se} \barrato{e} \barrato{solo} \barrato{se} \barrato{sono} \barrato{tutti} \barrato{P-dependent} \barrato{oppure}} {\color{red}$W^{22}$} {\color{red}\barrato{è} \barrato{regular.} \barrato{Le} \barrato{opinioni} \barrato{finali} \barrato{sono} \barrato{combinazioni} \barrato{convesse} \barrato{dei} \barrato{pregiudizi} \barrato{e} \barrato{delle} \barrato{opinioni} \barrato{finali} \barrato{dei} \barrato{P-independent.}}
    \item {\color{red}\barrato{Il} \barrato{caso} \barrato{che} \barrato{userò}} {\color{blue}A destra la condizione} di {\color{red}\barrato{più} \barrato{è} \barrato{il} \barrato{Corollario} \barrato{3.3:}} {\color{blue}stabilità asintotica:} il modello {\color{blue}FJ} è asintoticamente stabile se e solo se tutti gli agenti sono P-dependent. {\color{red}\barrato{Per}} {\color{blue}Succede per} esempio {\color{red}\barrato{quando} \barrato{la}} {\color{blue}su una} rete {\color{red}\barrato{è}} fortemente connessa {\color{red}\barrato{e}} {\color{blue}con} almeno un agente {\color{red}\barrato{è}} prejudiced. {\color{red}\barrato{E} \barrato{non} \barrato{serve} \barrato{nessuna} \barrato{aperiodicità.}}
    \item {\color{red}\barrato{L'intuizione:} \barrato{lo}} {\color{blue}In basso l'intuizione. Lo} scarto dall'equilibrio {\color{blue}$Vu$} evolve come {\color{red}$e(k+1)=\Lambda W\,e(k)$,} {\color{red}\barrato{senza}} {\color{red}$u$.} {\color{red}\barrato{Un} \barrato{agente} \barrato{prejudiced} \barrato{ne} \barrato{perde} \barrato{a} \barrato{ogni} \barrato{passo} \barrato{una} \barrato{frazione}} {\color{red}$1-\lambda_i$,} {\color{red}\barrato{quindi} \barrato{gli} \barrato{scarti} \barrato{si} \barrato{spengono} \barrato{solo}} {\color{blue}$e(k+1)=\Lambda W\,e(k)$: $u$ non compare. Quindi $\Lambda W$ decide} se {\color{red}\barrato{ogni} \barrato{agente} \barrato{è} \barrato{influenzato} \barrato{da} \barrato{un} \barrato{agente} \barrato{prejudiced.}} {\color{blue}le opinioni convergono, $u$ decide dove.}
    \item {\color{blue}Un agente prejudiced, a ogni passo, trattiene solo una frazione $\lambda_i$ degli scarti che ascolta: gli scarti si spengono solo se ogni agente è influenzato da un agente prejudiced.}
    \item {\color{blue}E non serve nessuna aperiodicità, a differenza del FDG.}
    \item {\color{blue}E se qualche agente è P-independent?}
    \item {\color{red}\barrato{Seconda} \barrato{domanda:} \barrato{dove} \barrato{convergono} \barrato{le} \barrato{opinioni?}}
  \end{discorso}
\end{slide}

\begin{slide}{m-fj-conv2}{Convergence with P-independent agents\nuova}{0:50}
  \begin{discorso}
    \item {\color{blue}Gli agenti P-independent non ascoltano nessun agente P-dependent. Mettendo prima i P-dependent, la matrice $W$ si divide in blocchi, e il blocco in basso a sinistra è zero.}
    \item {\color{blue}Le opinioni dei P-dependent dipendono dai prejudice e dalle opinioni dei P-independent; quelle dei P-independent evolvono da sole.}
    \item {\color{blue}Il blocco $W^{22}$ dice come i P-independent si ascoltano tra loro: è un modello FDG a sé, che nessun prejudice raggiunge.}
    \item {\color{blue}Il teorema generale di convergenza, a destra, dice due cose. La parte P-dependent è sempre asintoticamente stabile. E il modello FJ converge se e solo se tutti sono P-dependent, oppure converge il modello FDG dei P-independent.}
    \item {\color{blue}Quando converge, le opinioni finali dei P-dependent sono combinazioni convesse dei prejudice e delle opinioni finali dei P-independent.}
    \item {\color{blue}Quindi l'aperiodicità torna a servire, ma solo per i P-independent: lo vedremo in un esempio.}
    \item {\color{blue}Seconda domanda: dove convergono le opinioni?}
  \end{discorso}
\end{slide}

\begin{slide}{m-fj-final}{Final opinions and the augmented graph}{1:00}
  \begin{discorso}
    \item Con tutti gli agenti P-dependent le opinioni finali sono $x(\infty)=Vu$.
    \item Tre conseguenze. Le opinioni iniziali sono dimenticate: $x(0)$ non compare. Contano solo i {\color{red}\barrato{pregiudizi}} {\color{blue}prejudice} degli agenti prejudiced: se $\lambda_j=1$, la colonna $j$ di $V$ è nulla.
    \item E, per un dato $u$, il gruppo raggiunge un {\color{red}\barrato{consenso}} {\color{blue}consensus} se e solo se tutti gli agenti prejudiced hanno lo stesso {\color{red}\barrato{pregiudizio.}} {\color{blue}prejudice.}
    \item Il tutorial dà anche una seconda lettura, il modello {\color{red}\barrato{aumentato:}} {\color{blue}augmented:} a ogni agente $i$ si affianca un agente virtuale stubborn $i'$, che vale sempre $u_i$.
    \item La matrice $\bar A$ è row-stochastic, quindi {\color{blue}il modello augmented} è un modello FDG su $2n$ agenti. Nella figura l'agente 1 è l'unico prejudiced: a sinistra il modello FJ, a destra il grafo {\color{red}\barrato{aumentato}} {\color{blue}augmented} con l'agente virtuale $1'$.
    \item Ogni agente virtuale è un sink aperiodico; con tutti gli agenti P-dependent sono gli unici sink. {\color{red}\barrato{Letto} \barrato{con} \barrato{la} \barrato{Rosetta} \barrato{stone,}} {\color{blue}Più sink, tutti aperiodici: per il teorema del disagreement} è {\color{red}\barrato{il} \barrato{caso} \barrato{(iv):} \barrato{il} \barrato{disaccordo} \barrato{del} \barrato{Teorema} \barrato{2.4,}} {\color{blue}un disagreement,} con limite proprio $Vu$.
    \item In termini sociali: un {\color{red}\barrato{pregiudizio}} {\color{blue}prejudice} si comporta come un agente in più che non cambia mai opinione, messo accanto all'agente prejudiced.
    \item {\color{blue}Questo, letto al contrario, risponde alla domanda lasciata aperta.}
    \item {\color{red}\barrato{Terza} \barrato{domanda:} \barrato{quanto} \barrato{pesa} \barrato{ciascun} \barrato{agente?}}
  \end{discorso}
\end{slide}

\begin{slidedue}{m-back-question}{Back to the question\spostata}{1:15 {\color{red}\barrato{1:00}}}
  \begin{discorso}
    \item {\color{red}\barrato{Torno} \barrato{all'anello} \barrato{con} \barrato{gli} \barrato{agenti} \barrato{1} \barrato{e} \barrato{4} \barrato{stubborn.}}
    \item {\color{red}\barrato{Nel} \barrato{report} \barrato{c'è} \barrato{una} \barrato{seconda} \barrato{riscrittura,} \barrato{nella} \barrato{Figura} \barrato{8,} \barrato{lì}} {\color{blue}Riprendo la domanda lasciata aperta} con il {\color{red}\barrato{solo} \barrato{agente} \barrato{1:} \barrato{si} \barrato{tolgono}} {\color{blue}modello FDG. Sull'anello} gli agenti {\color{blue}1 e 4 sono} stubborn, {\color{red}\barrato{e} \barrato{le} \barrato{loro}} {\color{blue}non cambiano mai opinione: sono membri del gruppo, o} opinioni {\color{red}\barrato{diventano} \barrato{i} \barrato{pregiudizi} \barrato{degli} \barrato{agenti}} {\color{blue}fisse esterne} che {\color{red}\barrato{li} \barrato{ascoltavano.}} {\color{blue}gli altri ascoltano?}
    \item {\color{blue}Il modello augmented ha appena trasformato un prejudice in un agente stubborn virtuale. Letto al contrario: un agente stubborn è un prejudice per gli agenti che lo ascoltano.}
    \item {\color{red}\barrato{Per} \barrato{esempio} \barrato{l'agente}} {\color{blue}Lo vedo sull'anello: tolgo gli agenti 1 e 4. L'agente} 2 aggiorna come $x_2(k+1)=\tfrac12x_3(k)+\tfrac12x_1(0)$: è il modello FJ con $\lambda_2=\tfrac12$ e {\color{blue}prejudice} $u_2=x_1(0)$.
    \item Allo stesso modo {\color{red}$u_6=x_1(0)$} {\color{blue}l'agente 6 ha prejudice $x_1(0)$,} e {\color{red}$u_3=u_5=x_4(0)$,} {\color{blue}gli agenti 3 e 5 hanno prejudice $x_4(0)$,} tutti con $\lambda_i=\tfrac12$. Le traiettorie sono le {\color{red}\barrato{stesse.}} {\color{blue}stesse del FDG.}
    \item {\color{blue}\clic} Quindi la risposta è: valgono entrambe le letture.
    \item Come {\color{red}\barrato{membro}} {\color{blue}membri} del gruppo, {\color{red}\barrato{è} \barrato{la} \barrato{riscrittura} \barrato{del} \barrato{tutorial} \barrato{della} \barrato{slide} \barrato{precedente:}} {\color{blue}1 e 4 sono agenti FJ con} $\lambda_i=0$ e {\color{red}$u_i=x_i(0)$.} {\color{red}\barrato{Come}} {\color{blue}prejudice uguale alla loro} opinione {\color{red}\barrato{fissa} \barrato{esterna,} \barrato{è} \barrato{il} \barrato{pregiudizio} \barrato{degli} \barrato{agenti} \barrato{che}} {\color{blue}iniziale, cioè stubborn, e gli altri seguono} la {\color{red}\barrato{ascoltano,} \barrato{come} \barrato{qui.}} {\color{blue}regola del FDG.}
    \item {\color{blue}Come opinioni fisse esterne, sono il prejudice degli agenti che li ascoltano, come qui.}
    \item Le due letture danno le stesse traiettorie. {\color{red}\barrato{E} \barrato{il} \barrato{disaccordo}} {\color{blue}Il disagreement} dei due agenti stubborn è il {\color{red}\barrato{disaccordo}} {\color{blue}disagreement} di agenti ancorati a {\color{red}\barrato{pregiudizi}} {\color{blue}prejudice} diversi.
    \item {\color{blue}Stesse traiettorie: ma è lo stesso sistema? Lo vedremo negli esempi. Prima, la terza domanda: quanto pesa ciascun agente?}
    \item {\color{red}\barrato{Ora} \barrato{gli} \barrato{esempi} \barrato{del} \barrato{modello} \barrato{FJ,} \barrato{sulle} \barrato{stesse} \barrato{due} \barrato{reti.}}
  \end{discorso}
\end{slidedue}

\begin{slide}{m-fj-centrality}{Influence centrality}{0:40 {\color{red}\barrato{0:45}}}
  \begin{discorso}
    \item La social power richiede il {\color{red}\barrato{consenso,}} {\color{blue}consensus,} che nel modello FJ non è una proprietà del modello: dipende {\color{red}\barrato{da}} {\color{red}$u$.} {\color{blue}dai prejudice.} E il modello {\color{red}\barrato{aumentato}} {\color{blue}augmented} ha almeno $n$ sink, quindi per $n\ge2$ non va mai al {\color{red}\barrato{consenso.}} {\color{blue}consensus.}
    \item {\color{blue}Serve un'altra misura.} Come nel tutorial, {\color{red}\barrato{allora,}} prendo il modello FJ asintoticamente stabile {\color{red}\barrato{con}} {\color{blue}e} $u=x(0)$, l'ipotesi tradizionale della teoria {\color{red}\barrato{FJ.} \barrato{Così}} {\color{blue}FJ: è questa che dà un ruolo alle opinioni iniziali,} $x(\infty)=Vx(0)$.
    \item La influence centrality $c_i$ è il peso medio dell'opinione iniziale dell'agente $i$ nelle opinioni finali del gruppo: {\color{red}$c=n^{-1}V^\top\uno$,} {\color{red}\barrato{e}} la media {\color{red}\barrato{delle} \barrato{opinioni} \barrato{finali} \barrato{è}} {\color{red}$c^\top x(0)$.} {\color{blue}della colonna $i$ di $V$.}
    \item $V$ è stocastica, quindi $c$ è non negativo e somma a uno, come {\color{red}$w$.} {\color{blue}la social power.}
    \item {\color{blue}Ora gli esempi, sulle stesse due reti del FDG.}
    \item {\color{red}\barrato{Con}} {\color{red}$\Lambda=\alpha I$:} {\color{red}\barrato{per}} {\color{red}$\alpha=0$} {\color{red}\barrato{si} \barrato{ha}} {\color{red}$V=I$} {\color{red}\barrato{e}} {\color{red}$c$} {\color{red}\barrato{uniforme.}}
    \item {\color{red}\barrato{Il} \barrato{Lemma} \barrato{3.5} \barrato{collega} \barrato{le} \barrato{due} \barrato{misure:} \barrato{se}} {\color{red}$W$} {\color{red}\barrato{è} \barrato{fully} \barrato{regular,}} {\color{red}$c_\alpha$} {\color{red}\barrato{tende} \barrato{alla} \barrato{social} \barrato{power} \barrato{del} \barrato{modello} \barrato{FDG} \barrato{con}} {\color{red}$A=W$} {\color{red}\barrato{quando}} {\color{red}$\alpha$} {\color{red}\barrato{tende} \barrato{a} \barrato{1.}}
    \item {\color{red}\barrato{Torno} \barrato{ora} \barrato{all'agente} \barrato{stubborn,} \barrato{visto} \barrato{nei} \barrato{due} \barrato{modelli.}}
  \end{discorso}
\end{slide}

\parte{Parte 4 — Il modello FJ, esempi}{7:20 {\color{red}\barrato{5:00}}}
\begin{slide}{m-fj-networks}{FDG and FJ on the same networks\nuova}{0:50}
  \begin{discorso}
    \item {\color{blue}Per gli esempi uso le stesse due reti, senza self-loop: $W$ è la matrice della rete originale. Stesse opinioni iniziali, e il prejudice di ogni agente è la sua opinione iniziale, come nel modello FJ originale.}
    \item {\color{blue}Confronto sempre i due modelli a parità di peso: nel FDG il peso va sull'opinione attuale, nel FJ sul prejudice.}
    \item {\color{blue}Prima riga, il caso estremo: l'agente 1 stubborn è $a_{11}=1$ nel FDG, $\lambda_1=0$ con $u_1=x_1(0)$ nel FJ.}
    \item {\color{blue}Seconda riga, un solo agente con peso $t$. Le due equazioni sono uguali tranne un termine: nel FDG l'agente 1 dà peso $t$ alla sua opinione attuale $x_1(k)$, nel FJ al suo prejudice $u_1$. Quindi $1-\lambda_1=t$.}
    \item {\color{blue}Terza riga, tutti gli agenti con lo stesso peso 0.2: nel FDG $A=0.2I+0.8W$, nel FJ $\Lambda=0.8I$.}
    \item {\color{blue}Nelle prime due righe gli altri agenti restano come nella rete originale.}
    \item {\color{blue}Stesso peso, oggetto diverso: l'opinione attuale, che si muove ancora verso gli altri, o il prejudice, che non si muove mai.}
    \item {\color{blue}Comincio dal caso estremo, l'agente stubborn.}
  \end{discorso}
\end{slide}

\begin{slidedue}{m-state-input}{Stubborn agent: state or input}{1:15}
  \begin{discorso}
    \item {\color{blue}Prima riga della tabella: l'agente 1 stubborn.} Il tutorial {\color{red}\barrato{mostra} \barrato{che}} {\color{blue}riscrive così} un modello FDG con agenti {\color{red}\barrato{stubborn} \barrato{si} \barrato{riscrive} \barrato{come} \barrato{un} \barrato{modello} \barrato{FJ:}} {\color{blue}stubborn:} $\lambda_i=0$ e $u_i=x_i(0)$ per gli {\color{red}\barrato{agenti}} stubborn, $\lambda_i=1$ per gli altri.
    \item A sinistra il FDG con $A(1)$, a destra il FJ con $\lambda_1=0$: le traiettorie sono le {\color{red}\barrato{stesse.}} {\color{blue}stesse, come per la domanda.}
    \item Ma la stabilità no. Il modello FDG converge ma non è mai asintoticamente stabile, perché {\color{red}\barrato{1} \barrato{è}} {\color{blue}una matrice row-stochastic ha} sempre {\color{red}\barrato{un} \barrato{autovalore} \barrato{di}} {\color{red}$A$;} {\color{blue}l'autovalore 1;} il modello FJ invece è asintoticamente stabile, {\color{red}\barrato{per} \barrato{il} \barrato{Corollario} \barrato{3.3.}} {\color{blue}perché tutti gli agenti sono P-dependent.}
    \item \clic Per vederlo perturbo l'opinione iniziale dell'agente 1: $x_1(0)=0.5$, con $u$ invariato.
    \item Nel FDG l'agente 1 tiene 0.5 e il gruppo lo segue: la perturbazione non viene mai dimenticata. Nel FJ l'agente 1 torna su $u_1=0$ al primo passo, e il limite è lo stesso, $Vu$.
    \item Sulle matrici: la riga dell'agente 1 è $e_1^\top$ in $A(1)$ ed è zero in $\Lambda W$, quindi l'autovalore 1 di $A(1)$ diventa 0 in $\Lambda W$.
    \item In termini sociali: nel FDG la stubbornness è una memoria nello stato, mai dimenticata; nel FJ è un ingresso esterno, e lo stato dimentica $x(0)$.
    \item {\color{blue}Ora la seconda riga: un solo agente con peso un mezzo.}
    \item {\color{red}\barrato{Adesso} \barrato{posso} \barrato{rispondere} \barrato{alla} \barrato{domanda} \barrato{lasciata} \barrato{aperta.}}
  \end{discorso}
\end{slidedue}

\begin{slide}{m-one-prej}{A single prejudice}{0:50}
  \begin{discorso}
    \item {\color{red}\barrato{Stesso} \barrato{peso} \barrato{un} \barrato{mezzo} \barrato{sull'agente} \barrato{1} \barrato{nei} \barrato{due} \barrato{modelli:}} {\color{blue}Seconda riga della tabella, con $t=\tfrac12$:} a sinistra {\color{red}\barrato{come} \barrato{self-weight,} \barrato{nel}} {\color{blue}il} FDG con {\color{red}$A(\tfrac12)$;} {\color{blue}$A(\tfrac12)$, peso sull'opinione attuale;} a destra {\color{red}\barrato{come} \barrato{attaccamento} \barrato{al} \barrato{pregiudizio,} \barrato{nel}} {\color{blue}il} FJ con {\color{red}$\lambda_1=\tfrac12$} {\color{red}\barrato{e}} {\color{red}$u=x(0)$.} {\color{blue}$\lambda_1=\tfrac12$, peso sul prejudice.}
    \item La rete è fortemente connessa, quindi tutti gli agenti sono P-dependent e il modello FJ è asintoticamente stabile, anche se $G$ è periodico e non ha self-loop.
    \item Nel FJ $V=\uno e_1^\top$: tutto il gruppo finisce sul {\color{red}\barrato{pregiudizio}} {\color{blue}prejudice} dell'agente 1, zero, per ogni {\color{red}\barrato{attaccamento}} {\color{blue}attachment} positivo, per quanto piccolo. Nel FDG invece il {\color{red}\barrato{consenso}} {\color{blue}consensus} è a $3/7$.
    \item Sulla misura del potere: $c=e_1$ per ogni $\lambda_1<1$, l'agente 1 prende tutta la influence centrality; nel FDG, a $t=\tfrac12$, la sua social power è $2/7$.
    \item È di nuovo la testardaggine che ribalta la dinamica: un solo agente ancorato, che influenza tutti, decide dove finisce tutto il gruppo, per quanto debole sia il suo {\color{red}\barrato{attaccamento.}} {\color{blue}attachment.}
    \item E se tutti gli agenti hanno un {\color{red}\barrato{pregiudizio?}} {\color{blue}prejudice?}
  \end{discorso}
\end{slide}

\begin{slide}{m-all-prej}{Every agent prejudiced}{0:45}
  \begin{discorso}
    \item {\color{blue}Terza riga della tabella: tutti gli agenti hanno lo stesso peso 0.2, nel FDG come self-weight, nel FJ come attachment. I prejudice invece sono tutti diversi: ognuno ha come prejudice la propria opinione iniziale.}
    \item {\color{red}\barrato{Questo} \barrato{è} \barrato{il} \barrato{confronto} \barrato{equo:} \barrato{lo} \barrato{stesso} \barrato{peso} \barrato{0.2} \barrato{per} \barrato{tutti} \barrato{gli} \barrato{agenti.} \barrato{A} \barrato{sinistra} \barrato{il} \barrato{FDG} \barrato{con}} {\color{red}$A=0.2I+0.8W$,} {\color{red}\barrato{a} \barrato{destra} \barrato{il} \barrato{FJ} \barrato{con}} {\color{red}$\Lambda=0.8I$} {\color{red}\barrato{e}} {\color{red}$u=x(0)$.}
    \item Nel FDG {\color{red}\barrato{la} \barrato{matrice}} {\color{blue}$A$} è primitiva, {\color{red}\barrato{per} \barrato{il} \barrato{Corollario} \barrato{1.12,}} {\color{blue}perché la rete è fortemente connessa e ha self-loop,} e doubly-stochastic: {\color{red}\barrato{consenso}} {\color{blue}consensus} sulla media, 0.5.
    \item Nel FJ tutti {\color{red}\barrato{gli} \barrato{agenti}} sono prejudiced, quindi il modello è asintoticamente stabile, e {\color{red}$V=0.2(I-0.8W)^{-1}$} {\color{blue}$V$} ha tutte le entrate {\color{red}\barrato{positive.}} {\color{blue}positive: ogni prejudice entra in ogni opinione finale.} I {\color{red}\barrato{pregiudizi}} {\color{blue}prejudice} sono {\color{red}\barrato{diversi,} \barrato{quindi} \barrato{disaccordo.}} {\color{blue}diversi: disagreement.}
    \item Su una rete fortemente connessa il FDG non va mai in {\color{red}\barrato{disaccordo:}} {\color{blue}disagreement:} nel FJ il {\color{red}\barrato{disaccordo}} {\color{blue}disagreement} viene dai {\color{red}\barrato{pregiudizi.}} {\color{blue}prejudice.}
    \item In termini sociali: l'inerzia alla fine cede, l'ancora no; ancore diverse tengono il gruppo diviso, quello che il tutorial chiama un cleavage di opinioni.
    \item La slide seguente mostra perché questo è un {\color{red}\barrato{disaccordo}} {\color{blue}disagreement} anche nel senso del corso.
  \end{discorso}
\end{slide}

\begin{slidedue}{m-six-sinks}{Six prejudices, six sinks}{1:15}
  \begin{discorso}
    \item Riscrivo il modello FJ della {\color{red}\barrato{Figura} \barrato{10}} {\color{blue}terza riga} come modello FDG su dodici {\color{red}\barrato{agenti.}} {\color{blue}agenti, il modello augmented.} Ogni agente reale dà peso totale 0.8 agli altri agenti reali e 0.2 al suo agente virtuale, che è stubborn sul {\color{red}\barrato{pregiudizio}} {\color{blue}prejudice} $u_i$.
    \item Gli agenti reali formano una sola componente fortemente connessa, quella ombreggiata, che non è un sink. Collassandola si ottiene il condensation digraph: i sei agenti virtuali sono sei sink, aperiodici per i self-loop.
    \item {\color{red}\barrato{Teorema} \barrato{2.4}} {\color{blue}Più sink, tutti aperiodici: per il teorema del disagreement il modello va in disagreement, come} con {\color{red}\barrato{sei} \barrato{sink:} \barrato{è} \barrato{il} \barrato{caso} \barrato{(iv),} \barrato{come} \barrato{nel} \barrato{Remark} \barrato{2.7.}} {\color{blue}più agenti stubborn.} Lo stesso vale sull'anello.
    \item Quindi $x(\infty)=Vu$, con $V$ a entrate positive: ogni opinione finale è una combinazione convessa di tutti e sei i {\color{red}\barrato{pregiudizi.}} {\color{blue}prejudice.}
    \item Anche {\color{red}\barrato{la} \barrato{Figura} \barrato{9} \barrato{ha}} {\color{blue}con un solo agente prejudiced c'erano} sei sink, ma solo $1'$ {\color{red}\barrato{è}} {\color{blue}era} ascoltato: togliendo gli altri cinque resta un solo sink, e il {\color{red}\barrato{consenso.}} {\color{blue}consensus.}
    \item Così si chiude il filo del {\color{red}\barrato{disaccordo}} {\color{blue}disagreement} del FDG: un agente stubborn dà il {\color{red}\barrato{consenso,} \barrato{nella} \barrato{Figura} \barrato{2} \barrato{con}} {\color{red}$t=1$;} {\color{blue}consensus;} due danno il {\color{red}\barrato{disaccordo,}} {\color{blue}disagreement,} l'esempio della domanda; sei agenti virtuali, tutti ascoltati, danno il {\color{red}\barrato{disaccordo}} {\color{blue}disagreement} di questa slide.
    \item \clic Una nota sulle traiettorie. La {\color{red}\barrato{Figura} \barrato{10}} {\color{blue}figura} mostra solo i sei agenti reali; nel modello {\color{red}\barrato{aumentato}} {\color{blue}augmented} le traiettorie sarebbero dodici, e quelle dei sei agenti virtuali sono rette orizzontali sui {\color{red}\barrato{pregiudizi.}} {\color{blue}prejudice.}
    \item Ogni agente reale parte sulla retta del suo agente virtuale e la lascia appena la media delle opinioni che ascolta è diversa: sul ciclo subito, sull'anello gli agenti da 2 a 5 più tardi, quando arriva il cambiamento degli agenti 1 e 6.
    \item Il peso 0.2 sul suo agente virtuale lo richiama sempre, e si ferma su $Vu$, mai sul proprio {\color{red}\barrato{pregiudizio.}} {\color{blue}prejudice.} Nel FDG lo stesso 0.2 sta sull'opinione attuale, e l'agente segue il gruppo fino al {\color{red}\barrato{consenso.}} {\color{blue}consensus.}
    \item Ultimo caso: dove nessun {\color{red}\barrato{pregiudizio}} {\color{blue}prejudice} arriva.
  \end{discorso}
\end{slidedue}

\begin{slidedue}{m-aux}{Where no prejudice arrives}{1:10}
  \begin{discorso}
    \item {\color{red}\barrato{Questa} \barrato{rete} \barrato{serve} \barrato{per} \barrato{il} \barrato{caso}} {\color{blue}Ultimo caso,} escluso finora: non tutti gli agenti sono P-dependent.
    \item {\color{red}\barrato{Gli}} {\color{blue}La rete: gli} agenti 1, 2 e 3 formano un gruppo chiuso, un sink del condensation digraph; 4, 5 e 6 sono una catena che lo ascolta. {\color{red}\barrato{Un}} {\color{blue}In ogni esempio un} solo agente è prejudiced, con $\lambda=0.5$; ombreggiati gli agenti P-dependent.
    \item {\color{red}\barrato{Con}} {\color{red}$W_p$} {\color{red}\barrato{il} \barrato{gruppo} \barrato{chiuso} \barrato{è} \barrato{un} \barrato{ciclo} \barrato{periodico;} \barrato{con}} {\color{red}$W_a$} {\color{red}\barrato{l'arco} \barrato{da} \barrato{2} \barrato{a} \barrato{1} \barrato{lo} \barrato{rende} \barrato{aperiodico.}}
    \item {\color{red}\clic} {\color{red}\barrato{Caso} \barrato{(a):} \barrato{se}} {\color{blue}Esempio 1:} il prejudiced è l'agente 2, dentro il gruppo chiuso, {\color{blue}che è un ciclo periodico. Tutti raggiungono l'agente 2, quindi} tutti sono P-dependent. {\color{red}\barrato{Il} \barrato{modello} \barrato{è} \barrato{asintoticamente} \barrato{stabile}} {\color{blue}Ci aspettiamo stabilità asintotica,} e tutti {\color{red}\barrato{finiscono} \barrato{su}} {\color{red}$u_2=0.5$,} {\color{blue}sul prejudice dell'agente 2,} anche se il ciclo è periodico.
    \item {\color{red}\barrato{Caso} \barrato{(b):} \barrato{se}} {\color{blue}Esempio 2:} il prejudiced è l'agente 5, {\color{red}\barrato{fuori,} \barrato{solo} \barrato{4} \barrato{e} \barrato{5} \barrato{sono} \barrato{P-dependent.}} {\color{blue}fuori dal gruppo chiuso.} Il gruppo chiuso è {\color{red}\barrato{periodico,}} {\color{red}$W^{22}$} {\color{red}\barrato{non}} {\color{blue}P-independent: un modello FDG a sé, ed} è {\color{red}\barrato{regular,} \barrato{e}} {\color{blue}periodico. Ci aspettiamo che} il modello non {\color{red}\barrato{converge,} \barrato{per} \barrato{il} \barrato{Teorema} \barrato{3.2.}} {\color{blue}converga.}
    \item {\color{blue}Esempio 3: come il secondo, più l'arco da 2 a 1, che rende il gruppo chiuso aperiodico. Ci aspettiamo convergenza, ma non stabilità asintotica: il limite deve dipendere dalle opinioni iniziali del gruppo chiuso. Per verificarlo, nel pannello (d) cambio solo $x_2(0)$.}
    \item {\color{blue}\clic Esempio 1, in alto a sinistra: tutti vanno a 0.5, il prejudice dell'agente 2.}
    \item {\color{blue}Esempio 2, in alto a destra: il gruppo chiuso oscilla, e il modello non converge.}
    \item {\color{red}\barrato{Caso} \barrato{(c):} \barrato{con} \barrato{l'arco} \barrato{da} \barrato{2} \barrato{a} \barrato{1,}} {\color{red}$W^{22}$} {\color{red}\barrato{è} \barrato{regular.} \barrato{Il} \barrato{modello} \barrato{converge,} \barrato{ma} \barrato{non} \barrato{è} \barrato{asintoticamente} \barrato{stabile:} \barrato{gli} \barrato{agenti} \barrato{1,} \barrato{2,} \barrato{3} \barrato{e} \barrato{6} \barrato{finiscono} \barrato{a} \barrato{0.4,} \barrato{gli} \barrato{agenti} \barrato{4} \barrato{e} \barrato{5} \barrato{a} \barrato{0.65.}}
    \item {\color{red}\barrato{Caso} \barrato{(d):} \barrato{se} \barrato{cambio} \barrato{solo}} {\color{red}$x_2(0)$,} {\color{blue}Esempio 3, in basso: converge; 1, 2, 3 e 6 vanno a 0.4, 4 e 5 a 0.65. Portando $x_2(0)$} da 0.5 a {\color{red}\barrato{1,} \barrato{i} \barrato{limiti}} {\color{blue}1} diventano 0.6 e {\color{red}\barrato{0.75.} \barrato{Il}} {\color{blue}0.75: il} limite dipende {\color{blue}davvero} dalle opinioni iniziali dei P-independent.
    \item In termini sociali: gli agenti che nessun {\color{red}\barrato{pregiudizio}} {\color{blue}prejudice} raggiunge formano un modello FDG a sé; hanno bisogno dell'aperiodicità e ricordano le loro opinioni iniziali.
    \item Infine, le due misure di potere.
  \end{discorso}
\end{slidedue}

\begin{slidedue}{m-lemma}{Centrality meets social power}{1:15 {\color{red}\barrato{1:00}}}
  \begin{discorso}
    \item {\color{blue}Ho due misure per la stessa domanda: quanto pesa l'opinione iniziale di ciascun agente.}
    \item {\color{blue}La social power, nel FDG: la quota di ciascun agente nell'opinione su cui il gruppo si mette d'accordo. Richiede il consensus.}
    \item {\color{blue}La influence centrality, nel FJ: quanto pesa in media l'opinione iniziale di ciascuno in opinioni finali che possono essere diverse. Richiede la stabilità asintotica e $u=x(0)$.}
    \item {\color{blue}Sono misure diverse: la social power pesa lo stato iniziale di un modello che ricorda $x(0)$, la influence centrality l'ingresso di un modello che lo dimentica. Per questo non le confronto direttamente, ma in parallelo.}
    \item {\color{blue}Si incontrano nel lemma a destra: se $W$ è fully regular, cioè il FDG con $A=W$ raggiunge il consensus, e tutti hanno la stessa susceptibility $\alpha$, per $\alpha$ che tende a 1 la influence centrality tende alla social power del FDG.}
    \item Per {\color{red}\barrato{confrontare} \barrato{le} \barrato{due} \barrato{misure}} {\color{blue}vederlo} prendo $W=A(\tfrac12)$, la rete con self-weight un mezzo sull'agente {\color{red}\barrato{1,} \barrato{e} \barrato{la} \barrato{stessa} \barrato{suscettibilità}} {\color{red}$\alpha$} {\color{red}\barrato{per} \barrato{tutti,} \barrato{con}} {\color{blue}1: è primitiva, quindi fully regular. E} $u=x(0)$.
    \item {\color{red}$W$} {\color{red}\barrato{è} \barrato{primitiva,} \barrato{per} \barrato{il} \barrato{Corollario} \barrato{1.12,} \barrato{quindi} \barrato{fully} \barrato{regular:} \barrato{il} \barrato{Lemma} \barrato{3.5} \barrato{si} \barrato{applica.}}
    \item {\color{blue}Cosa ci aspettiamo.} Per {\color{red}$\alpha=0$,} {\color{red}$V=I$} {\color{blue}$\alpha=0$ ogni agente tiene la sua opinione iniziale: $V=I$,} e la influence centrality è uniforme, $1/6$. {\color{red}\barrato{Al} \barrato{crescere} \barrato{di}} {\color{red}$\alpha$,} {\color{red}$c_1$} {\color{red}\barrato{sale} \barrato{fino} \barrato{alla} \barrato{social} \barrato{power}} {\color{red}$w_1=2/7$,} {\color{red}\barrato{e}} {\color{red}$c_2$} {\color{red}\barrato{scende} \barrato{fino}} {\color{blue}Per $\alpha$ che tende} a {\color{red}$w_2=1/7$.} {\color{red}\barrato{Sul} \barrato{ciclo}} {\color{red}$c_2$} {\color{red}\barrato{scende} \barrato{prima} \barrato{anche} \barrato{sotto}} {\color{red}$w_2$.} {\color{blue}1, $c$ tende a $w$: $2/7$ per l'agente 1, $1/7$ per gli altri.}
    \item {\color{blue}\clic Ed è quello che succede: $c_1$ sale da $1/6$ fino a $2/7$, $c_2$ scende fino a $1/7$. Sul ciclo $c_2$ prima scende anche sotto $1/7$.}
    \item In termini sociali: man mano che gli agenti si aprono agli altri, l'inerzia dell'agente 1 conta di più.
    \item La social power dice chi guida un gruppo che si mette d'accordo; la influence centrality dice chi pesa di più in un gruppo che resta in {\color{red}\barrato{disaccordo.}} {\color{blue}disagreement.} La seconda diventa la prima quando le ancore si allentano.
    \item Per concludere, i due modelli a confronto.
  \end{discorso}
\end{slidedue}

\parte{Parte 5 — Conclusioni}{1:30}
\begin{slide}{m-table}{FDG and FJ side by side}{0:40}
  \begin{discorso}
    \item La tabella mette a confronto i due modelli riga per riga.
    \item La resistenza: nel FDG è un self-weight, un'inerzia; nel FJ un {\color{red}\barrato{attaccamento,}} {\color{blue}attachment,} un'ancora.
    \item L'agente stubborn: nel FDG è una condizione sulla sola matrice, $a_{ii}=1$; nel FJ sulla matrice e sull'ingresso, $\lambda_i=0$ e $u_i=x_i(0)$. La sua opinione è uno stato nel FDG, un ingresso nel FJ.
    \item La stabilità asintotica: mai nel FDG; nel FJ se e solo se tutti gli agenti sono P-dependent.
    \item Su una rete fortemente connessa il FDG dà {\color{red}\barrato{consenso}} {\color{blue}consensus} o nessuna convergenza, mai {\color{red}\barrato{disaccordo;}} {\color{blue}disagreement;} il FJ dà {\color{red}\barrato{disaccordo}} {\color{blue}disagreement} ogni volta che i {\color{red}\barrato{pregiudizi}} {\color{blue}prejudice} che contano sono diversi.
    \item L'aperiodicità serve sempre nel FDG; nel FJ solo per gli agenti P-independent.
    \item E la misura del potere: la social power richiede il {\color{red}\barrato{consenso;}} {\color{blue}consensus;} la influence centrality richiede la stabilità asintotica e $u=x(0)$.
  \end{discorso}
\end{slide}

\begin{slide}{m-messages}{Take-home messages}{0:50}
  \begin{discorso}
    \item Tre messaggi da portare a casa.
    \item Primo, inerzia e ancora. Un self-weight rallenta un agente ma lo lascia seguire gli altri: su una rete fortemente connessa e periodica basta un self-weight fra 0 e 1 per trasformare le oscillazioni in {\color{red}\barrato{consenso.}} {\color{blue}consensus.} Un {\color{red}\barrato{pregiudizio}} {\color{blue}prejudice} non cede mai: un solo agente prejudiced che influenza tutti decide dove finisce il gruppo, e {\color{red}\barrato{pregiudizi}} {\color{blue}prejudice} diversi lo tengono diviso.
    \item Secondo, stato e ingresso. La stessa opinione stubborn è uno stato nel modello FDG, che non è mai asintoticamente stabile, e un ingresso nel modello FJ, asintoticamente stabile quando tutti gli agenti sono P-dependent. Il grafo {\color{red}\barrato{aumentato}} {\color{blue}augmented} riporta gli ingressi a stati.
    \item Terzo, il potere. La social power cresce gradualmente con il self-weight; un agente prejudiced che influenza tutti prende tutta la influence centrality; e con {\color{red}$\Lambda=\alpha I$} {\color{blue}la stessa susceptibility per tutti} e $W$ fully regular, la influence centrality tende alla social power quando $\alpha$ tende a 1.
  \end{discorso}
\end{slide}

\begin{slidepag}{\the\numexpr\pg{m-messages}+1\relax}{Thank you!}{}
  \begin{discorso}
    \item Grazie per l'attenzione. Sono a disposizione per le domande.
  \end{discorso}
\end{slidepag}

\end{document}
